\documentclass[11pt]{article}
% Shared print layout. The web's provisional design and MathJax are independent.
\usepackage[utf8]{inputenc}
\usepackage[T1]{fontenc}
\usepackage[spanish,es-nodecimaldot]{babel}
\usepackage{lmodern}
\usepackage{amsmath,amssymb}
\usepackage{graphicx}
\usepackage{geometry}
\geometry{a4paper,left=20mm,right=20mm,top=25mm,bottom=22mm,headheight=16pt,headsep=8mm}
\usepackage{fancyhdr,lastpage}
\usepackage{enumitem}
\usepackage{xcolor}
\usepackage{pgfplots}
\pgfplotsset{compat=1.18}
\usepgfplotslibrary{fillbetween}
\usepackage[hidelinks,pdfencoding=auto]{hyperref}
\hypersetup{pdflang=es-ES,pdfauthor={},pdfsubject={Matemáticas Aplicadas a las Ciencias Sociales II}}
\pagestyle{fancy}
\fancyhf{}
\fancyhead[L]{\small Matemáticas Aplicadas a las Ciencias Sociales II}
\fancyhead[R]{\small 2.\textsuperscript{o} Bachillerato}
\fancyfoot[L]{\footnotesize Programación lineal II | Edición revisada 1.0}
\fancyfoot[R]{\footnotesize Página \thepage\ de \pageref{LastPage}}
\renewcommand{\headrulewidth}{0.4pt}
\setlength{\parindent}{0pt}
\setlength{\parskip}{0.5em}
\setlist{itemsep=0.3em,topsep=0.3em}
\newcommand{\apartado}[1]{\subsection*{#1}}
\pgfplotsset{
  revision/.style={
    width=0.94\linewidth,height=7.5cm,
    axis lines=middle,grid=major,grid style={gray!18},
    tick label style={font=\small},label style={font=\small},
    xlabel={$x$},ylabel={$y$},samples=160,
    legend style={at={(0.5,-0.16)},anchor=north,draw=none,font=\small},
    clip=true
  }
}

\hypersetup{pdftitle={Programación lineal II - Soluciones revisadas}}
\begin{document}
\section*{Ejercicio 1: baterías y beneficio máximo}

Sean $x$ las baterías A e $y$ las baterías B. El modelo es
\[
\max Z=130x+140y,\qquad
\begin{cases}x+y\geq10,\\y-x\leq10,\\3x+2y\leq120,\\x\geq0,\quad y\geq0.\end{cases}
\]
La tercera restricción procede de dividir $150x+100y\leq6\,000$ entre 50. Los beneficios del enunciado ya son beneficios unitarios: no se restan de nuevo los gastos de producción.

Las intersecciones factibles y sus valores son
\begin{center}
\begin{tabular}{c|r}
Vértice & $Z$ (euros)\\\hline
$(0,10)$ & $1\,400$\\
$(10,0)$ & $1\,300$\\
$(40,0)$ & $5\,200$\\
$(20,30)$ & $6\,800$
\end{tabular}
\end{center}
La intersección de $y=x+10$ y $3x+2y=120$ da $5x=100$, luego $(x,y)=(20,30)$.

\textbf{Certificado para toda la región:}
\[Z=54(3x+2y)+32(y-x)\leq54\cdot120+32\cdot10=6\,800.\]
En $(20,30)$ se alcanza la igualdad y se cumplen $x+y=50\geq10$, $y-x=10$ y $150x+100y=6\,000$. Así,
\[\boxed{20\text{ baterías A},\quad30\text{ baterías B},\quad Z_{\max}=6\,800\text{ euros}}.\]
El óptimo es entero y único, pues la igualdad exige simultáneamente las dos rectas anteriores.
\begin{center}
\begin{tikzpicture}
\begin{axis}[revision,height=5.2cm,xmin=-2,xmax=44,ymin=-2,ymax=36,xtick={0,10,20,30,40},ytick={0,10,20,30},xlabel={Baterías A ($x$)},ylabel={Baterías B ($y$)}]
\addplot[fill=gray!20,draw=black,thick] coordinates {(0,10) (10,0) (40,0) (20,30)} -- cycle;
\addplot[only marks,mark=*,mark size=2pt] coordinates {(0,10) (10,0) (40,0) (20,30)};
\node[anchor=south] at (axis cs:20,30) {\small óptimo $(20,30)$};
\end{axis}
\end{tikzpicture}
\end{center}
\clearpage
\section*{Ejercicio 2: misiones, solo formulación}
Este ejercicio pide \textbf{formular, sin resolver}. Sean $x$ las misiones de Observación de la Tierra e $y$ los programas de Transporte Espacial.
\subsection*{Función objetivo}
\[\max F(x,y)=0{,}6x+0{,}4y=\frac35x+\frac25y.\]
La función objetivo es la que da el enunciado; no representa por sí misma una cantidad de euros.
\subsection*{Restricciones del presupuesto y la inversión}
El gasto total no puede superar 2,4 millones:
\[200\,000x+100\,000y\leq2\,400\,000\quad\Longleftrightarrow\quad2x+y\leq24.\]
Debe \textbf{superar}, y no igualar, 2 millones:
\[200\,000x+100\,000y>2\,000\,000\quad\Longleftrightarrow\quad2x+y>20.\]
\subsection*{Restricciones de las cantidades}
Al menos cuatro misiones: $x\geq4$. El número de misiones no supera la mitad de los programas: $x\leq y/2$, equivalente a $y\geq2x$.
\subsection*{Modelo completo}
\[
\boxed{\begin{gathered}
\max\ F(x,y)=\frac35x+\frac25y\\
2x+y\leq24,\qquad2x+y>20,\\
x\geq4,\qquad y\geq2x,\\
x,y\in\mathbb Z_{\geq0}.
\end{gathered}}
\]
Para el modelo continuo utilizado en representación gráfica, se relaja la integralidad a cantidades reales no negativas. Las decisiones finales de misiones y programas son enteras. En ese caso $2x+y>20$ equivale a $2x+y\geq21$.

La formulación termina aquí: no se calculan vértices ni óptimos, conforme a lo solicitado en el examen.
\clearpage
\section*{Ejercicio 3: recinto y mínimo}

La región está definida por
\[x+4y\geq5,\quad x+2y\geq4,\quad7x+5y\leq35,\quad x\geq0.\]
Al intersectar rectas frontera y descartar puntos que incumplen alguna restricción, obtenemos:
\begin{center}
\begin{tabular}{c|c|c}
Vértice & Fronteras que se cortan & $F=2x+y$\\\hline
$(0,2)$ & $x=0$, $x+2y=4$ & $2$\\
$(3,1/2)$ & $x+4y=5$, $x+2y=4$ & $13/2$\\
$(5,0)$ & $x+4y=5$, $7x+5y=35$ & $10$\\
$(0,7)$ & $x=0$, $7x+5y=35$ & $7$
\end{tabular}
\end{center}
Por ejemplo, restando $x+2y=4$ de $x+4y=5$, resulta $y=1/2$ y después $x=3$. Todos los vértices indicados cumplen las cuatro restricciones.

\textbf{Certificado del mínimo para toda la región:}
\[F=\frac12(x+2y)+\frac32x\geq\frac12\cdot4+0=2.\]
La igualdad exige $x=0$ y $x+2y=4$, es decir, $(0,2)$, que es factible. Por tanto,
\[\boxed{F_{\min}=2\quad\text{únicamente en }(0,2)}.\]
\begin{center}
\begin{tikzpicture}
\begin{axis}[revision,height=5.8cm,xmin=-0.5,xmax=5.8,ymin=-0.5,ymax=7.8,xtick={0,1,2,3,4,5},ytick={0,2,4,6,7}]
\addplot[fill=gray!20,draw=black,thick] coordinates {(0,2) (3,0.5) (5,0) (0,7)} -- cycle;
\addplot[densely dotted,domain=0:5] {(5-x)/4};
\addplot[dashed,domain=0:5] {(4-x)/2};
\addplot[domain=0:5] {(35-7*x)/5};
\addplot[only marks,mark=*,mark size=1.7pt] coordinates {(0,2) (3,0.5) (5,0) (0,7)};
\node[anchor=west] at (axis cs:0.15,2.1) {\small mínimo $(0,2)$};
\end{axis}
\end{tikzpicture}
\end{center}
El contorno del polígono está incluido. No se añade $y\geq0$ como hipótesis: se deduce de las restricciones del recinto.
\clearpage
\section*{Anexo documental: frutos rojos}
\textbf{Material conservado de F020. No forma parte del examen F022, no añade puntos y no sustituye a su ejercicio 2.}
\subsection*{Enunciado de la variante}
Una frutería vende surtidos A y B. Cada A contiene 75 g de arándanos y 100 g de frambuesas, y se vende a 2,40 euros. Cada B contiene 75 g de arándanos y 50 g de frambuesas, y se vende a 1,80 euros. Se dispone de 3,75 kg de arándanos y 4 kg de frambuesas. La cantidad de surtidos A vendidos es menor o igual al doble de la de B.

El original pide formular, \textbf{sin resolver}, un problema para maximizar el ``beneficio''.
\subsection*{Qué se puede formular con los datos disponibles}
Sean $x$ los surtidos A e $y$ los B. Convertimos a gramos: $3{,}75$ kg son $3\,750$ g y $4$ kg son $4\,000$ g.
\begin{align*}
75x+75y\leq3\,750&\quad\Longleftrightarrow\quad x+y\leq50,\\
100x+50y\leq4\,000&\quad\Longleftrightarrow\quad2x+y\leq80,\\
x&\leq2y.
\end{align*}
La función calculable con precios de venta es la de \textbf{ingresos}:
\[
\boxed{\begin{gathered}
\max\ I(x,y)=2{,}40x+1{,}80y\\
x+y\leq50,\qquad2x+y\leq80,\\
x\leq2y,\qquad x,y\in\mathbb Z_{\geq0}.
\end{gathered}}
\]
El modelo continuo sustituye la integralidad por $x,y\geq0$ reales.
\subsection*{Mejora pendiente del enunciado original}
Los precios de venta no son beneficios netos. Para maximizar el beneficio neto faltan costes o márgenes unitarios, o una hipótesis explícita que permita ignorar costes variables. No se inventan esos datos. Otra corrección posible sería pedir \textbf{ingresos máximos} en el enunciado.

No se resuelve el modelo ni se incorpora esta variante a la puntuación del examen.
\end{document}
